\chapter*{Цели и задачи}
\addcontentsline{toc}{chapter}{Введение}
\label{ch:intro_lab9}

\textbf{Цель лабораторной работы:} научиться производить кластерный анализ данных на основе метода K-средних.

\textbf{Основные задачи:}
\begin{itemize}
    \item получение навыков рефакторинга кода в проектах машинного обучения;
    \item изучение принципов определения оптимального количества кластеров в методах кластерного анализа;
    \item изучение возможностей языка Python для реализации кластерного анализа;
    \item освоение метрик оценки качества кластеризации (Silhouette Score, Calinski-Harabasz Index, Davies-Bouldin Index);
    \item визуализация результатов кластеризации с использованием PCA и дендрограмм.
\end{itemize}

\textbf{Теоретическое обоснование:}

Кластеризация — это разбиение множества объектов на подмножества (кластеры) по заданному критерию. Каждый кластер включает максимально схожие между собой объекты. В отличие от классификации, кластеризация не требует размеченных данных, что делает её методом обучения без учителя.

В данной лабораторной работе используется метод \textbf{K-средних (K-Means)} — один из наиболее популярных алгоритмов кластеризации. Алгоритм работает следующим образом:

\begin{enumerate}
    \item Выбирается количество кластеров $K$
    \item Инициализируются $K$ центроидов (случайным образом или методом K-Means++)
    \item Каждый объект относится к ближайшему центроиду
    \item Пересчитываются центроиды как средние значения объектов в кластере
    \item Шаги 3-4 повторяются до сходимости
\end{enumerate}

Для оценки качества кластеризации используются следующие метрики:

\begin{itemize}
    \item \textbf{Silhouette Score} — показывает, насколько объект похож на свой кластер по сравнению с другими кластерами (значения от -1 до 1)
    \item \textbf{Calinski-Harabasz Index} — отношение межкластерной дисперсии к внутрикластерной (чем выше, тем лучше)
    \item \textbf{Davies-Bouldin Index} — среднее сходство между кластерами (чем ниже, тем лучше)
\end{itemize}

В качестве набора данных для выполнения лабораторной работы используется набор данных \textbf{Abalone}, который содержит информацию о морских моллюсках. Данные включают 8 признаков (длина, диаметр, высота, весовые характеристики и пол) и целевую переменную Rings (возраст в кольцах).

\endinput